
\documentclass{article}
\usepackage{lgd}


% Attempt to make hyperref and algorithmic work together better:
\newcommand{\theHalgorithm}{\arabic{algorithm}}
\newcommand{\Bern}{{\mbox{Bern}}}
\newcommand{\Lap}{{\mbox{Lap}}}
% \newcommand{\bX}{{\bf X}}
\newcommand{\bXt}{\tilde{\bf X}}
\newcommand{\bXh}{\hat{\bf X}}
\newcommand{\bNsup}[1]{{\bf N}^{(#1)}}
\newcommand{\bNtsup}[1]{\tilde{\bf N}^{(#1)}}
\newcommand{\bZsup}[1]{{\bf Z}^{(#1)}}
\newcommand{\bZtsup}[1]{\tilde{\bf Z}^{(#1)}}
\newcommand{\bXhsup}[1]{\hat{\bf X}^{(#1)}}
\newcommand{\bXtsup}[1]{\tilde{\bf X}^{(#1)}}
\newcommand{\bzerod}{{\bf 0}_d}
\newcommand{\boned}{{\bf 1}_d}

\newcommand{\SpLap}{{\mbox{SpLap}}}
\newcommand{\sgn}{{\mbox{sign}}}
\newcommand{\iid}{\stackrel{iid}{\sim}}
\newcommand{\deq}{\stackrel{(d)}{=}}

\newcommand{\equalto}[2]{\underset{\scriptstyle\overset{\mkern4mu\verteq}{#2}}{#1}}
\newcommand{\verteq}{\rotatebox{90}{$\,=$}}



%% some caligraphic symbols
\newcommand{\cA}{{\cal A}}
\newcommand{\cB}{{\cal B}}
\newcommand{\cC}{{\cal C}}
\newcommand{\cD}{{\cal D}}
\newcommand{\cE}{{\cal E}}
\newcommand{\cF}{{\cal F}}
\newcommand{\cG}{{\cal G}}
\newcommand{\cH}{{\cal H}}
\newcommand{\cI}{{\cal I}}
\newcommand{\cJ}{{\cal J}}
\newcommand{\cK}{{\cal K}}
\newcommand{\cL}{{\cal L}}
\newcommand{\cM}{{\cal M}}
\newcommand{\cN}{{\cal N}}
\newcommand{\cO}{{\cal O}}
\newcommand{\cP}{{\cal P}}
\newcommand{\cQ}{{\cal Q}}
\newcommand{\cR}{{\cal R}}
\newcommand{\cS}{{\cal S}}
\newcommand{\cT}{{\cal T}}
\newcommand{\cU}{{\cal U}}
\newcommand{\cV}{{\cal V}}
\newcommand{\cW}{{\cal W}}
\newcommand{\cX}{{\cal X}}
\newcommand{\cY}{{\cal Y}}
\newcommand{\cZ}{{\cal Z}}

\DeclareMathAlphabet{\pazocal}{OMS}{zplm}{m}{n}
\newcommand{\calD}{{\pazocal {D}}}
\newcommand{\calG}{{\pazocal {G}}}


%% some bold symbols
\newcommand{\ba}{{\bf a}}
\newcommand{\bb}{{\bf b}}
\newcommand{\bc}{{\bf c}}
\newcommand{\bd}{{\bf d}}
\newcommand{\be}{{\bf e}}
%\newcommand{\bf}{{\bf f}}
\newcommand{\bg}{{\bf g}}
\newcommand{\bh}{{\bf h}}
\newcommand{\bj}{{\bf j}}
\newcommand{\bk}{{\bf k}}
\newcommand{\bl}{{\bf l}}
\newcommand{\bm}{{\bf m}}
\newcommand{\bn}{{\bf n}}
\newcommand{\bo}{{\bf o}}
\newcommand{\bp}{{\bf p}}
\newcommand{\bq}{{\bf q}}
\newcommand{\br}{{\bf r}}
\newcommand{\bs}{{\bf s}}
\newcommand{\bt}{{\bf t}}
\newcommand{\bu}{{\bf u}}
\newcommand{\bv}{{\bf v}}
\newcommand{\bw}{{\bf w}}
\newcommand{\bx}{{\bf x}}
\newcommand{\by}{{\bf y}}
\newcommand{\bz}{{\bf z}}

\newcommand{\bA}{{\bf A}}
\newcommand{\bB}{{\bf B}}
\newcommand{\bC}{{\bf C}}
\newcommand{\bD}{{\bf D}}
\newcommand{\bE}{{\bf E}}
\newcommand{\bF}{{\bf F}}
\newcommand{\bG}{{\bf G}}
\newcommand{\bH}{{\bf H}}
\newcommand{\bI}{{\bf I}}
\newcommand{\bJ}{{\bf J}}
\newcommand{\bK}{{\bf K}}
\newcommand{\bL}{{\bf L}}
\newcommand{\bM}{{\bf M}}
\newcommand{\bN}{{\bf N}}
\newcommand{\bO}{{\bf O}}
\newcommand{\bP}{{\bf P}}
\newcommand{\bQ}{{\bf Q}}
\newcommand{\bR}{{\bf R}}
\newcommand{\bS}{{\bf S}}
\newcommand{\bT}{{\bf T}}
\newcommand{\bU}{{\bf U}}
\newcommand{\bV}{{\bf V}}
\newcommand{\bW}{{\bf W}}
\newcommand{\bX}{{\bf X}}
\newcommand{\bY}{{\bf Y}}
\newcommand{\bZ}{{\bf Z}}


\newcommand{\tA}{\tilde{A}}
\newcommand{\tB}{\tilde{B}}
\newcommand{\tC}{\tilde{C}}
\newcommand{\tD}{\tilde{D}}
\newcommand{\tE}{\tilde{E}}
\newcommand{\tF}{\tilde{F}}
\newcommand{\tG}{\tilde{G}}
\newcommand{\tH}{\tilde{H}}
\newcommand{\tI}{\tilde{I}}
\newcommand{\tJ}{\tilde{J}}
\newcommand{\tK}{\tilde{K}}
\newcommand{\tL}{\tilde{L}}
\newcommand{\tM}{\tilde{M}}
\newcommand{\tN}{\tilde{N}}
\newcommand{\tO}{\tilde{O}}
\newcommand{\tP}{\tilde{P}}
\newcommand{\tQ}{\tilde{Q}}
\newcommand{\tR}{\tilde{R}}
\newcommand{\tS}{\tilde{S}}
\newcommand{\tT}{\tilde{T}}
\newcommand{\tU}{\tilde{U}}
\newcommand{\tV}{\tilde{V}}
\newcommand{\tX}{\tilde{X}}
\newcommand{\tY}{\tilde{Y}}
\newcommand{\tZ}{\tilde{Z}}

\newcommand{\tu}{\tilde{u}}
\newcommand{\tx}{\tilde{x}}
\newcommand{\ty}{\tilde{y}}
\newcommand{\tgamma}{\tilde{\gamma}}
%\newcommand{\ta}{\tilde{a}}

\newcommand{\field}[1]{\mathbb{#1}}
\newcommand{\fC}{\field{C}} 
\newcommand{\fE}{\field{E}} 
\newcommand{\fN}{\field{N}} % natural numbers
\newcommand{\fR}{\field{R}} % real numbers
\newcommand{\fS}{\field{S}}
\newcommand{\fZ}{\field{Z}} % integers

\newcommand{\uh}{\hat{u}}
\newcommand{\vh}{\hat{v}}
\newcommand{\xh}{\hat{x}}
\newcommand{\yh}{\hat{y}}
\newcommand{\zh}{\hat{z}}
\newcommand{\Uh}{\hat{U}}
\newcommand{\Xh}{\hat{X}}
\newcommand{\Yh}{\hat{Y}}
\newcommand{\Zh}{\hat{Z}}
\newcommand{\cXh}{\hat{\cX}}
\newcommand{\bUh}{\hat{\bU}}
\newcommand{\cYh}{\hat{\cY}}
\newcommand{\cZh}{\hat{\cZ}}

\newcommand{\bQstar}{\bQ^{\star}}
\newcommand{\Astar}{A^{\star}}
\newcommand{\Bstar}{B^{\star}}
\newcommand{\Qstar}{Q^{\star}}

\newcommand{\tsnr}{{\text{snr}}}
\newcommand{\tvar}{{\text{Var}}}
\newcommand{\cmle}{\text{\rm cmle}}
\newcommand{\cmse}{\text{\rm cmse}}
\newcommand{\mmse}{\text{\rm mmse}}
\newcommand{\mmle}{\text{\rm mmle}}
\newcommand{\minimax}{\text{\rm minimax}}
\newcommand{\mse}{\text{\rm mse}}
\newcommand{\mle}{\text{\rm mle}}

\newcommand{\diag}{\mbox{diag}}
\newcommand{\JSD}{\mbox{JSD}}
\newcommand{\conv}{\text{\rm conv}}

\newcommand{\mathify}[1]{\ifmmode{#1}\else\mbox{$#1$}\fi}
\newcommand{\Expec}[1]{\mathify{\mathbb{E}\left[#1\right]}}
\newcommand{\Var}[1]{\mathify{{\rm VAR}\left[#1\right]}}
\newcommand{\Cov}{\mathify{{\rm Cov}}}

% \newcommand{\Esubs}[2]{\mathify{\mathbb{E}_{#1}\left[#2\right]}}

\newcommand{\Err}{\mathrm{Err}}


%% bold symbols and some other symbols for ratless paper
\newcommand{\mathUO}{{\rm Unif}(\mathcal{O}(n))}
\newcommand{\Xb}[1]{{\bf X^{(#1)}}}
\newcommand{\Bb}[1]{{\bf B^{(#1)}}}
\newcommand{\Xe}[1]{ X^{(#1)}}
\newcommand{\Be}[1]{ B^{(#1)}}
\newcommand{\Zb}[1]{{\bf Z^{(#1)}}}
\newcommand{\Ze}[1]{ Z^{(#1)}}
\newcommand{\Yb}[1]{{\bf Y^{(#1)}}}
\newcommand{\Ye}[1]{ Y^{(#1)}}
\newcommand{\Ub}[1]{{\bf U^{(#1)}}}
\newcommand{\Ue}[1]{ U^{(#1)}}
\newcommand{\Vb}[1]{{\bf V^{(#1)}}}
\newcommand{\Ve}[1]{ V^{(#1)}}
\newcommand{\Wb}[1]{{\bf W^{(#1)}}}
\newcommand{\We}[1]{ W^{(#1)}}
\newcommand{\Wbtilde}[1]{{\bf \widetilde{W}^{(#1)}}}
\newcommand{\Wetilde}[1]{ \widetilde{W}^{(#1)}}
\newcommand{\Ah}{\hat{A}}
\newcommand{\Ab}{{\bf A}}

\newcommand{\Qb}{{\bf Q}}
\newcommand{\Xbhat}{{\bf \hat{X}}}
\newcommand{\Xbh}[1]{{\bf \hat{X}^{(#1)}}}
\newcommand{\mb}[1]{{\bf m^{(#1)}}}
\newcommand{\me}[1]{ m^{(#1)}}
\newcommand{\mbhat}[1]{{\bf \hat{m}^{(#1)}}}
\newcommand{\mehat}[1]{{\hat{m}^{(#1)}}}

\usepackage{tikz}
\usepackage{pstricks}
\usetikzlibrary{arrows,shapes}
\usetikzlibrary{fit}

%% math operators
\newcommand\argmin{\mathop{\mbox{{\rm argmin}}}\limits}
\newcommand\argmax{\mathop{\mbox{{\rm argmax}}}\limits}
\newcommand\maximize{\mathop{\mbox{{\rm maximize}}}\limits}
\newcommand\minimize{\mathop{\mbox{{\rm minimize}}}\limits}
\newcommand{\bea}{\begin{eqnarray}}
\newcommand{\eea}{\end{eqnarray}}
\newcommand{\beas}{\begin{eqnarray*}}
\newcommand{\eeas}{\end{eqnarray*}}
\usepackage{amsfonts}
\newcommand{\Espace}{{\tilde{\mathcal{F}}[0,T]}}
\newcommand{\suchthat}{\;\ifnum\currentgrouptype=16 \middle\fi|\;}
\newcommand{\dbar}{\;\ifnum\currentgrouptype=16 \middle \fi\|\;}


\allowdisplaybreaks

% if you use cleveref..
\usepackage[capitalize,noabbrev]{cleveref}


% Todonotes is useful during development; simply uncomment the next line
%    and comment out the line below the next line to turn off comments
%\usepackage[disable,textsize=tiny]{todonotes}
\usepackage[textsize=tiny]{todonotes}




\newcommand{\noa}[1]{{\color{red}\textbf{AN}: #1}}
\newcommand{\berivan}[1]{{\color{orange}\textbf{BI}: #1}}
\newcommand{\wn}[1]{{\color{green}\textbf{WN}: #1}}


%%%%%%%%%%%%%%%%%%% ADDED BY WEI-NING %%%%%%%%%%%%%%%
\newcommand{\lp}{\left(}
\newcommand{\rp}{\right)}
\newcommand{\lb}{\left[}
\newcommand{\rb}{\right]}
\newcommand{\lbp}{\left\{}
\newcommand{\rbp}{\right\}}
\newcommand{\lba}{\left\lvert}
\newcommand{\rba}{\right\rvert}
\newcommand{\lV}{\left\lVert}
\newcommand{\rV}{\right\rVert}
\newcommand{\mv}{\middle\vert}
\newcommand{\lan}{\langle}
\newcommand{\ran}{\rangle}
\newcommand{\ra}{\rightarrow}
\newcommand{\eqDef}{\coloneqq}
\newcommand{\mcal}{\mathcal}

\newcommand{\diid}{\overset{\text{i.i.d.}}{\sim}}
%%%%%%%%%%%%%%%%%%%%%%%%%%%%%%%%%%%%%%%%%%%%%%%%%%%%%

\newcommand{\punif}{P_{\text{unif}}}

\begin{document}
\lecturetitle{3}{Information Theory}{Albert No}{03/18/2024}


%%% Cross Etnropy Loss
%%%%%%%%%%%%%%%%%%%%%%%%%%%%%%%%%%%%%%%%%%%%%%
\section{Cross Entropy Loss} \label{sec:intro-cross}
\begin{figure}[h]
    \centering
    \includegraphics[width=0.3\textwidth]{images/classification.png}
    \caption{Classifier $f$}
\end{figure}

Consider a classification problem where data input $x\in\cX$ is from the set of possible inputs (such as $\fR^n$)
and the corresponding label is $y\in [M] = \{1, 2, \dots, M\}$.
A classifier $f:\cX \rightarrow \cM$ produces the probability distribution $f(x) = (f(x)_1, \dots, f(x)_M)$ 
where $f(i)$ indicates the ``the probability of $x$ corresponds to label $i$.''
Note that $\cM$ is a set of probability measures on set of labels.
The cross entropy is defined by
\begin{align}
    \ell(f(x), y) = \log \frac{1}{f(x)_y}
\end{align}


In one hand, it is weird looking function:
1) it is not clear why we are using $\log$ function,
2) it is not clear why we are using reciprocal.
Of course, the definition make sense, since ``if your model assigns the higher value to the true label, the loss value is lower.''
However, such property holds for any decreasing function, it does not have to be $\log 1/f(y)$.

Some of you may already know that the cross entropy can be derived from the Kullback-Leibler (KL) divergence.
For two probability distributions $P, Q\in \cM$, 
KL divergence is defined by (we will cover KL divergence more thoroughly later)
\begin{align*}
    D(P\| Q) = \sum_{i} P(i) \log \frac{P(i)}{Q(i)}
\end{align*}

It measures a ``distance'' between two probability distribution\footnote{strictly speaking
it is not a distance since it is not symmetric.}.
Let $y^{(o)}\in\fR^M$ is an one-hot vector represents ground truth, i.e.,
\begin{align}
    y^{(o)}(i) = \begin{cases}
        1 &\mbox{ if $y=i$}\\
        0 &\mbox{ otherwise}
    \end{cases}    
\end{align}
The one-hot vector $y^{(o)}$ can be viewed as a ground truth distribution of input $x$,
which is deterministically pointing the true label $y$.
Then, it is natural to consider the ``distance'' between the true distribution $y^{(o)}$ and the output of classifier $f(x)$
to measure the performance.

The KL divergence between $f(x)$ and $y^{(o)}$ is
\begin{align}
    D(y^{(o)}\| f(x)) 
    =&  \sum_{i=1}^M y^{(o)}_i \log \frac{f(x)_i}{y^{(o)}_i}\\
    =& \log \frac{1}{f(x)_y}\\ 
    =& \ell(f(x),y).
\end{align}
In other words, the KL divergence is equivalent to cross entropy,
which justifies the usage of the cross entropy loss function.

However, the concern remains.
The KL divergence is also weird looking.
Why are we measuring ``distance'' between distributions in such way?
Is there another way of measuring ``distance''?
Why is it called divergence? instead of ``distance''?

To answer these questions, we need to understand the basic information measures in information theory
from entropy, mutual information, to KL divergence (or often called relative entropy).


\begin{remark}
    We will see later, KL divergence plays a crucial role in Generative Adversarial Networks (GAN).
\end{remark}

\subsection{Notation} \label{subsec:notation}
We denote upper case $X$ to indicate the (discrete) random variable,
where its support is finite $X\in\{1, \dots, M\}$.
For its realization, we use lower case $x$.
The probability mass function is $P_X(x) = \Pr{X= x}$.
In most cases, the logarithm function is base 2 otherwise specified.

%%% Entropy
%%%%%%%%%%%%%%%%%%%%%%%%%%%%%%%%%%%%%%%%%%%%%%
\section{Entropy} \label{sec:entropy}
For discrete random variable $X\in\{1, \dots, M\}$,
we define suprisal.
\begin{definition}
    Surprisal of $X$ at $x$ is given by
    \begin{align}
        S(x) = \log \frac{1}{P_X(x)}.
    \end{align}
\end{definition}
This is natural definition since if rare events happens (where $P_X(x)$ is small), then it is ``surprising''.

The entropy would be an expected suprisal of random variable $X$.
\begin{definition}
    Entropy of random variable $X$ is given by
    \begin{align}
        H(X) =& \Expec{S(X)}\\
        =& \Expec{-\log P_X(X)}\\
        =& \sum_{x=1}^M P_X(x) \log \frac{1}{P_X(x)}.
    \end{align}
\end{definition}

\begin{example}
    Let $X\in\{1, 2\}$ be a binary random variable where $P_X(1) = 0.01$ and $P_X(2) = 0.99$, then
    \begin{align}
        H(X) = 0.01 \log \frac{1}{0.01} + 0.99 \log \frac{1}{0.99}.
    \end{align}
\end{example}

%%% Properties of Entropy
\subsection{Properties of Entropy}
\begin{theorem}[Property of Entropy 1]
    Entropy is non-negative, and it is 0 if and only if $X$ is deterministic.
\end{theorem}
\begin{proof}
    Since the probability is always $0\leq P_X(x)\leq 1$, we have
    \begin{align}
        H(X) =& \Expec{\log \frac{1}{P_{X}(X)}}\\
        \geq & \Expec{\log 1}\\
        =& 0.
    \end{align}
    The equality holds if and only if $P_X(x)\log P_X(x) = 0$ for all $x$,
    which is true only if $X$ is deterministic.
    Note that $0\log 0 = 0$.
\end{proof}

\begin{theorem}[Jensen's Inequality]
Let $f$ be a concave function, then
\begin{align}
    \Expec{f(X)} \leq f(\Expec{X}).
\end{align}
\end{theorem}
Since the definition of concave function is that
\begin{align}
    \lambda f(x) + (1-\lambda) f(y) \leq f(\lambda x + (1-\lambda) y)
\end{align}
for all $x, y$ and $0\leq \lambda \leq 1$,
in other words, weighted average of $f$ is smaller or equal to $f$ at weighted averaged input.
Then, Jensen's inequality is a generalized version of it,
where $\Expec{f(X)}$ is weighted average of $f$ and $f(\Expec{X})$ is a function value at weighted average input.
For strictly concave function $f$, Jensen's inequality holds with equality if and only if $f(x)$ is constant.

\begin{theorem}[Property of Entropy 2]
    The maximum entropy is $\log M$.
\end{theorem}
\begin{proof}
Since $\log$ is a concave function, by Jensen's inequality
\begin{align}
    H(X) =& \Expec{\log \frac{1}{P_X(X)}}\\
    \leq& \log \Expec{\frac{1}{P_X(X)}}\\
    =& \log \sum_{x=1}^M P_X(x) \frac{1}{P_{X}(x)}\\
    =& \log M.
\end{align}
\end{proof}

Equality of the above inequality holds if and only if $\frac{1}{P_X(x)}$ is constant,
which implies $X$ is uniform on $\{1,\dots,M\}$, i.e., $P_X(x) = 1/M$ for all $x$.

\begin{remark}
    Entropy can also be viewed as an ``amount of information''.
    It is also natural since surprising (rare) event contains more information.
    This will become even more clear when we discuss lossless compression.
\end{remark}
\begin{definition}[Joint Entropy]
    For jointly distributed random variables $X_1$ and $X_2$ with joint probability mass function $P_{X_1, X_2}(x_1, x_2)$.
    The joint entropy is given by
    \begin{align}
        H(X_1, X_2) = \sum_{x_1, x_2} P_{X_1, X_2} (x_1, x_2) \frac{1}{P_{X_1, X_2} (x_1, x_2)}.
    \end{align}
\end{definition}
Note that $X_1$ and $X_2$ may have different support, i.e., $X_1\in[M] = \{1, \dots,M\}$ and $X_2 \in [N] = \{1, \dots, N\}$.


\begin{theorem}[Property of Entropy 3]
    If $X_1$ and $X_2$ are independent,
    \begin{align}
        H(X_1, X_2) = H(X_1) + H(X_2).
    \end{align}
\end{theorem}
\begin{proof}
If $X_1$ and $X_2$ are independent, then $P_{X_1, X_2} (x_1, x_2) = P_{X_1}(x_1) P_{X_2}(x_2)$ for all $x_1, x_2$.
\begin{align}
    H(X_1, X_2) =& \sum_{x_1, x_2} P_{X_1, X_2} (x_1, x_2) \log \frac{1}{P_{X_1}(x_1)P_{X_2}(x_2)}\\
    =& \sum_{x_1, x_2} P_{X_1, X_2} (x_1, x_2) \log \frac{1}{P_{X_1}(x_1)}
    + \sum_{x_1, x_2} P_{X_1, X_2} (x_1, x_2) \log \frac{1}{P_{X_2}(x_2)}\\
    =& \sum_{x_1}  \log \frac{1}{P_{X_1}(x_1)}\sum_{x_2}P_{X_1, X_2} (x_1, x_2)
    + \sum_{x_2} \log \frac{1}{P_{X_2}(x_2)}\sum_{x_1} P_{X_1, X_2} (x_1, x_2) \\
    =& \sum_{x_1}  \log \frac{1}{P_{X_1}(x_1)} P_{X_1} (x_1)
    + \sum_{x_2} \log \frac{1}{P_{X_2}(x_2)} P_{X_2} (x_2) \\
    =& H(X_1)+H(X_2).
\end{align}
\end{proof}
It is clear that if $X_1$ and $X_2$ are highly correlated,
the amount of information that we obtain from $(X_1, X_2)$ should be similar to
the amount of information from $X_1$.

\begin{exercise}
    If $H(X_1, X_2)=H(X_1) +H(X_2)$, then does this imply independence?
\end{exercise}

The Property 3 can be generalized to multi variables.
If $X_1, \dots, X_n$ are independent, i.e., $P_{X_1,\dots, X_n}(x_1,\dots, x_n) = \prod_{i=1}^n P_{X_i}(x_i)$ ,then
\begin{align}
    H(X_1,\dots, X_n) = \sum_{i=1}^n H(X_i).
\end{align}

\begin{theorem}[Property of Entropy 4]
For ``wrong'' probability mass function $Q(x)$,
\begin{align}
    H(X) \leq \Expec{\log \frac{1}{Q(X)}} = \sum_{x} P_X(x) \log \frac{1}{Q(x)}.
\end{align}
\end{theorem}
\begin{example}
    Let $X$ be the same binary distribution $P_X(1) = 0.01$ and $P_X(2) = 0.99$ where $Q(1) = 0.1$ and $Q(2) = 0.9$, then the right hand side would be
    \begin{align}
        \Expec{\log \frac{1}{Q(X)}} = 0.01 \log \frac{1}{0.1} + 0.99 \log\frac{1}{0.9}.
    \end{align}
\end{example}
\begin{proof}
    Since $\log$ is concave function, by Jensen's inequality~\cite{jensen1906fonctions},
    \begin{align}
        H(X) - \Expec{\log \frac{1}{Q(X)}} =& \Expec{\log \frac{Q(X)}{P_X(X)}}\\
        \leq& \log \Expec{\frac{Q(X)}{P_X(X)}}\\
        =& \log \sum_x P_X(x)\frac{Q(x)}{P_X(x)}\\
        =& \log \sum_x Q(x)\\
        =& \log 1 =0.
    \end{align}
    Equality holds if and only if $\frac{Q(x)}{P_X(x)}$ is constant, which implies $Q = P_X$.
\end{proof}

\begin{remark}
    From the above proof, we have
    \begin{align}
        \Expec{\log \frac{P_X(X)}{Q(X)}} = \sum_{x} P_X(x) \log\frac{P_X(x)}{Q(x)}\geq 0
    \end{align}
    which implies
    \begin{align}
        D(P_X\|Q) \geq 0.
    \end{align}
    In other words, the KL divergence is always non-negative and 0 if and only if $P_X=Q$.
\end{remark}

\begin{remark}
    The above proof is indeed incomplete since there is an edge case where $P_X(x) = 0$ for some $x$.
    The rigorous proof is
    \begin{align}
        H(X) - \Expec{\log \frac{1}{Q(X)}} =& \sum_{x}P_X(x)\log \frac{Q(x)}{P_X(x)}\\
        =& \sum_{x: P_X(x) > 0} P_X(x)\log \frac{Q(x)}{P_X(x)}\\
        \leq& \log \sum_{x:P_X(x) > 0} P_X(x)\frac{Q(x)}{P_X(x)}\\
        =& \log \sum_{x:P_X(x) > 0} Q(x)\\
        \leq & \log 1 =0.
    \end{align}
    To achieve equality, both the first (Jensen) and the second ($\sum_{x:P_X(x)>0} Q(x) \leq 1$) inequality should hold with equality.
    In other words, 1) $Q(x)/P_X(x)=c$ is constant for all $x$ and 2) $Q(x) = 0$ if $P_X(x) = 0$.
    This can only be true if $Q = P_X$.
\end{remark}

\begin{definition}
    The right hand side of the above property is called ``cross entropy,''
    where we happen to use wrong pmf $Q(x)$ instead of true pmf $P_X(x)$.
\end{definition}

\subsection{Cross Entropy Loss Revisited}
Recall that
\begin{align}
    KL(P_X\|Q) =& \sum_{x} P_X(x) \log \frac{P_X(x)}{Q(x)}\\
    =& -\sum_{x} P_X(x) \log \frac{1}{P_X(x)} + \sum_x \log\frac{1}{Q(x)}\\
    =& -H(P_X) +\sum_x \log P_X(x)\frac{1}{Q(x)}
\end{align}
where the second term is called cross entropy which coincides with the cross entropy loss
\begin{align}
    CE(y^{(o)}\|f(x)) = \sum_m y^{(o)}_m \log \frac{1}{f(x)_m}.
\end{align}
This is the reason why we call it ``cross entropy loss.''
Thus, even when the true label is ``soft'' distribution (instead of one-hot), we have
\begin{align}
    KL(y^{(soft)} \| f(x)) = -H( y^{(soft)}) + CE(y^{(soft)} \| f(x)).
\end{align}
Since the true label is given and fixed, minimizing KL divergence is equivalent to minimizing cross entropy loss.


\begin{remark}
    We often abuse the notation $H(X) = H(P_X)$.
\end{remark}


%%% Mutual Information
%%%%%%%%%%%%%%%%%%%%%%%%%%%%%%%%%%%%%%%%%%%%%%%%%%%%%%%%%%%%%%%%%%
\section{Mutual Information}
\begin{definition}
   Let $X$ and $Y$ be jointly distributed random variables, then the conditional entropy is given by
   \begin{align}
       H(X|Y) =& \Expec{\log\frac{1}{P_{X|Y}(X|Y)}}\\
       =& \sum_{x, y} P_{X, Y}(x, y) \log\frac{1}{P_{X|Y}(x|y)}\\
       =& \sum_{y} P_{Y}(y)\sum_x P_{X|Y}(x|y) \log\frac{1}{P_{X|Y}(x|y)}\\
       =& \sum_{y} P_{Y}(y) H(X|Y=y)
   \end{align}
\end{definition}
Note that $H(X|Y=y)$ is remaining uncertainty of $X$ when $Y$ is known to be $y$.
On the other hand, $H(X|Y)$ is (expected) remaining uncertainty of $X$ when $Y$ is provided (thus average over $H(X|Y=y)$).
Intuitively, some $y$ might be informative while some $y$ are not.

\begin{definition}
    Mutual information between $X$ and $Y$ is given by
    \begin{align}
        I(X;Y) = H(X) - H(X|Y).
    \end{align}
\end{definition}
Mutual information is difference between ``uncertainty of $X$'' and ``remaining uncertainty of $X$ given $Y$,''
and therefore it implies the amount information that $Y$ providing about $X$.

\begin{lemma}
    Mutual information is symmetric
    \begin{align}
        I(X;Y) = I(Y;X)
    \end{align}
\end{lemma}

Also, we can re-write the mutual information using KL divergence
\begin{align}
    I(X;Y) = D(P_{X, Y} \| P_X P_Y)
\end{align}
\begin{proof}
\begin{align}
    I(X;Y) =& \Expec{\frac{1}{P_{X}(X)} - \frac{1}{P_{X|Y}{X|Y}}}\\
=& \Expec{\frac{P_{X|Y}(X|Y)}{P_{X}(X)}}\\
=& \Expec{\frac{P_{X,Y}(X,Y)}{P_{X}(X)P_{Y}(Y)}}\\
=& D(P_{X,Y}\|P_X P_Y).
\end{align}
\end{proof}
This implies that the mutual information is a ``distance'' between joint distribution and product of marginal (as if $X$ and $Y$ are independent) distributions.
Thus, it is also a measure of ``how $X$ and $Y$ are correlated.''


\bibliographystyle{abbrv}
\bibliography{scribe-ref}
\end{document}